\documentclass[14pt, russian]{article}
\usepackage{graphicx}
\linespread{1.3}
\usepackage{amsbsy}
\usepackage[T2A]{fontenc}
\usepackage[cp1251]{inputenc}
\usepackage[russian]{babel}
\usepackage{epstopdf}
\usepackage[autostyle]{csquotes}
\usepackage[a4paper,left=2.5cm,right=2.5cm,top=2.5cm,bottom=2.5cm]{geometry}%
\usepackage{cuted}   
\usepackage{lipsum}  
\usepackage[intoc,nocfg,russian]{nomencl}
\newcommand{\nomencl}[2]{#1 --- #2\nomenclature{#1}{#2}}
\setlength{\nomlabelwidth}{3em}
\setlength{\nomitemsep}{-\parsep}
\setcounter{tocdepth}{2}
\renewcommand{\nomlabel}[1]{#1 ---}
\makenomenclature
\usepackage{wrapfig}
\let\vec=\mathbf
\def\mprp{\mbox{\tiny $\bot$}}
\def\mprl{\mbox{\tiny $\|$}}
\def\th{\mbox{th}}
\def\D{\mathrm{d}}
\def\beq{\begin{eqnarray}}
\def\eeq{\end{eqnarray}}
\def\ee{\varepsilon}
\def\lm{\lambda}
\newcommand{\ii}{\mathrm{i}} % мнимая единица
\newcommand{\dd}{\mathrm{d}} % дифференциал
\newcommand{\eee}{\mathrm{e}} % основание натуральных логарифмов

\newcommand{\bs}{\boldsymbol}

\begin{document}
\begin{center}
Ответ на отзыв рецензента на статью Д. М. Шленева и др.\\
«Перенос излучения в равновесной электронной плазме
с учетом резонансного комптоновского процесса»

\vspace{5mm}
Уважаемая редакция
\end{center}

В ответ на замечание рецензента нами были внесены в статью следующие изменения:
\begin{enumerate}
	\item На с. 3 в конце левой колонке исправлено выражение $\zeta=(M_1^2+m^2)/eB=2+2B_e/B$.
	\item На 2 с. в правой колонке первого абзаца исправлено выражение $O(\alpha^2)$. Явно указано, что $\alpha$ -- постоянная тонкой структуры.
	\item На 2 с. левой колонки было исправлено условие, при котором рассеяние фотонов на ионах можно пренебречь: "При этом мы не учитываем взаимодействие фотонов с ионами, поскольку из-за их высокой массы сечение рассеяния на них мало по сравнению с рассеянием на электронах при энергиях фотонов, намного превосходящих энергию ионного циклотронного резонанса.".
	\item В формуле (2) было добавлено определение энергий $E, E'$ начального и конечного электронов. В формулах для $f_E$ и $f_{E'}$ действительно в общем случае должен присутствовать химический потенциал электронов, однако в списке цитированной литературы зачастую рассматривают модели с нулевым химическим потенциалом. В связи с этим в статье явно указано, что мы рассматриваем химический потенциал равный нулю.
	Однако учет ненулевого химического не~меняет структуру уравнений в данной статье -- он будет входить в функции распределения электронов, которые содержатся в концентрации электронов $n_e$.
	\item На с.4 добавлен переход от формулы (17) к формуле (18) согласно нумерации исправленной статьи, он отражен в формуле (15). Был исправлен коэффициент для функции распределения фотона -- он в действительности отличается от коэффициента в спектре. Также был исправлен множитель для энергии функции распределения фотона (18) (эта ошибка не повлияла на результаты, так как строились исходя из конечной формы спектра (17)).
	\item Во Введении уточнено определение относительно слабого магнитного поля. Оно определяется по отношению к критическому значению. Также в аннотации и заключение указано точное значение, относительно которого поле будем считать относительно слабым. Также было добавлено явное выражение для критического значения $B_e = m^2/e$
	\item После формулы (14) было явно дано определение комптоновской длины волны $\lambda_C$.
	\item Ссылка для определения нормальных мод была выбрана исходя из работы Муштукова, потому что она послужила отправной точкой для наших исследований. Тем не менее мы полностью согласны с отзывом, что будет уместно сослаться на фундаментальные работы, поэтому была добавлена ссылка на монографию Гинзбурга В. Л. Распространение электромагнитных волн в плазме.
	\item Были исправлены указанные искажения в транскрибации и написании фамилий авторов: 
			\begin{itemize}
				\item Нишимура заменено на Нисимура согласно системе Поливанова;
				\item Изменено Зварт (S.P. Zwart) на Портегис Цварт (S. Portegies Zwart);
				\item Изменено Мезарос на Мешарос.
			\end{itemize}
	\item Исправлены указанные в отзыве мелкие недочеты.
	
\end{enumerate} 

\vspace{5mm}
С уважением,

Д. М. Шленев, Д. А. Румянцев, А. А. Ярков.
%
%Функция
%%
%\begin{equation}
%\label{f}
%f_{\varepsilon} (x) = \frac{1}{\pi} \frac{\varepsilon}{x^2 + \varepsilon^2} \; 
%(\varepsilon >0)
%\end{equation}
%%
%образует дельта-образную последовательность при $\varepsilon\to 0$ [1]':
%%
%\begin{equation}
%\lim_{\varepsilon\to 0} \frac{1}{\pi} \frac{\varepsilon}{x^2 + \varepsilon^2} 
%= \delta(x).
%\end{equation}
%%
%Поскольку в формуле (21) данной статьи ширина резонанса имеет конечное 
%значение, определяемое параметром $\varepsilon = E_n''\Gamma^{s''}_n/2$, 
%квадрат модуля пропагатора может быть преобразован к виду, подобному 
%выражению~(\ref{f}), но это подобие понимается в смысле обобщенных функций, т. 
%е. в смысле сходимости к дельта-функции при интегрировании:
%\begin{eqnarray}
%	I(\varepsilon) = \pi \varphi(0) + \frac{\varepsilon}{4} P.V. 
%\int_{-\infty}^{\infty} \frac{\varphi(x)}{x^2} dx- \frac{\pi\varepsilon^2}{2} 
%\left.\frac{\partial^2 \varphi(x)}{\partial x^2}\right|_{x=0} + 
%O(\varepsilon^3),
%\end{eqnarray}
%где интеграл вычисляется в смысле главного значения (P.V. -- Principal Value), 
%$\Phi(s)$ -- Фурье-образ функции $\varphi(x)$.
%
%Таким образом замена
%\begin{eqnarray}
%	&&\left|\frac{1}{(p+q)^2_{\mprl}-M_n^2-\ii E_n''\Gamma^{s''}_n/2}\right|^2 
%= 
%	\frac{1}{((p+q)^2_{\mprl}-M_n^2)^2 + (E_n''\Gamma^{s''}_n/2)^2} = 
%	\\
%	\nonumber
%	&&=\frac{1}{E_n''\Gamma^{s''}_n/2}
%	\frac{E_n''\Gamma^{s''}_n/2}{(p+q)^2_{\mprl}-M_n^2-\ii 
%E_n''\Gamma^{s''}_n/2} 
%	\simeq\frac{2\pi}{E_n''\Gamma_n^{s''}} \delta((p+q)^2_{\mprl}-M_n^2)
%\end{eqnarray}
%с точностью до $E_n''\Gamma^{s''}_n\ll \left|(p+q)^2_{\mprl}-M_n^2\right|$ 
%имеет смысл только в подынтегральном выражении. 
%
%Реальная часть в точке резонанса равна нулю $(p+q)^2_{\mprl}-M_n^2=0$, что 
%соответствует тому, что виртуальный электрон приобретает массу $M_n$ и 
%становится реальным.
%
%В связи с этим было добавлено приложение 1, поясняющее выражение (3), 
%используя преобразование Фурье, а также уточнён текст, который указывает на 
%последующее интегрирование формулы (21), что и позволяет сделать эффективную 
%замену на дельта-функцию.
%
%1. Гельфанд, И.М. Обобщенные функции и действия над ними/ И. М. Гельфанд, Г. 
%Е. Шилов. – Москва.:Государственное издательство физико-математической 
%литературы: 1958. – 470 с.
%


\end{document}